\problem{mixed\_precision\_accumulation}{Mixed-Precision Accumulation}{1}

Run the assignment's repeated-addition example with FP32 accumulation, FP16 accumulation, and FP16 values accumulated in FP32. Compare the resulting values and explain their accuracy.

\deliverable{A 2--3 sentence response.}
\answer

\problem{benchmarking\_mixed\_precision}{Benchmarking Mixed Precision}{2}

\partquestion{a}{For the specified ToyModel under FP16 autocast, identify the dtype of the parameters, first linear output, layer-normalization output, logits, loss, and gradients.}
\deliverable{The dtype of each listed component.}

\begin{table}[htbp]
\centering\small
\begin{tabular}{@{}ll@{}}
\toprule
Component & Dtype \\
\midrule
Model parameters & \\
First feed-forward layer output & \\
Layer-normalization output & \\
Predicted logits & \\
Loss & \\
Gradients & \\
\bottomrule
\end{tabular}
\caption{Dtypes in the FP16 autocast ToyModel.}
\label{tab:autocast-dtypes}
\end{table}

\partquestion{b}{Which layer-normalization operations are sensitive to mixed precision? Does BF16 require the same special treatment as FP16, and why?}
\deliverable{A 2--3 sentence response.}
\answer

\partquestion{c}{Add optional BF16 mixed precision to the benchmark and compare forward and backward timings against full precision for every model size. Describe trends with model size.}
\deliverable{A 2--3 sentence response with timings and commentary.}
\answer
